\input{preamble/preamble.tex}

\geometry{margin=0.85in}
\AtBeginDocument{\small}

\newcommand{\ExamNameField}{\noindent\textbf{Name:}\ \rule{0.7\linewidth}{0.4pt}\par\medskip}

\begin{document}
	\ExamTitleBlock{11th grade}{Term 2 Catch-Up: C1 Sample Mean and Sample Standard Deviation (Solutions)}{}

	\ExamSection{Problems}
	\begin{ExamProblems}
		\item
		\subsection*{Problem 1}
		\subsection*{Problem description}
		A group of Grade 11 students manages a small digital-product studio that sells logo templates and social media packs. They tracked one random sample of \textbf{14 weekly net payouts} (USD) during a semester pilot to understand whether their business cash flow is stable enough to support reinvestment decisions.

		\textbf{Sample data (USD):}
		\[118,\ 124,\ 131,\ 127,\ 136,\ 142,\ 125,\ 139,\ 133,\ 121,\ 145,\ 129,\ 137,\ 126\]

		Compute $\bar{x}$, $s^2$, and $s$. Then write a short interpretation of what the variability says about weekly income risk for this studio.

		\subsection*{C1}
		We compute the sample mean and the sample standard deviation for the 14 weekly payouts.

		Step 1: Compute $\bar{x}$.
		\[
		\bar{x} = \frac{\sum x_i}{n}
		= \frac{118+124+131+127+136+142+125+139+133+121+145+129+137+126}{14}
		= \frac{1833}{14}
		= 130.93\ \text{USD}
		\]

		Step 2: Compute $s$.
		\[
		s = \sqrt{\frac{\sum (x_i-\bar{x})^2}{n-1}}
		\]
		Using $\bar{x}=130.93$:
		\[
		\sum (x_i-\bar{x})^2 = 844.93
		\]
		\[
		s = \sqrt{\frac{844.93}{14-1}} = \sqrt{\frac{844.93}{13}} = \sqrt{64.99} = 8.06\ \text{USD}
		\]

		The average weekly payout is about 130.93 USD, and a typical week is about 8.06 USD away from this average, indicating moderate week-to-week income risk.

		\newpage
		\item
		\subsection*{Problem 2}
		\subsection*{Problem description}
		A student entrepreneur receives income from two channels: short-form video monetization and freelance thumbnail design. She collected two independent samples of weekly net income (USD) to decide which channel is more predictable before choosing where to invest more time next month.

		\textbf{Sample A (Video monetization, USD):}
		\[62,\ 75,\ 68,\ 91,\ 57,\ 84,\ 73,\ 95,\ 66,\ 88\]

		\textbf{Sample B (Freelance design, USD):}
		\[54,\ 58,\ 61,\ 63,\ 57,\ 60,\ 56,\ 65,\ 59,\ 62,\ 58\]

		For each sample, compute the mean and sample standard deviation. Based on your results, identify which channel appears more volatile and explain why this matters for planning reliable monthly cash flow.

		\subsection*{C1}
		We compute center and spread for each channel, then compare volatility.

		\textbf{Sample A}

		Step 1: Compute $\bar{x}_A$.
		\[
		\bar{x}_A = \frac{\sum x_i}{n_A}
		= \frac{62+75+68+91+57+84+73+95+66+88}{10}
		= \frac{759}{10}
		= 75.90\ \text{USD}
		\]

		Step 2: Compute $s_A$.
		\[
		s_A = \sqrt{\frac{\sum (x_i-\bar{x}_A)^2}{n_A-1}}
		\]
		Using $\bar{x}_A=75.90$:
		\[
		\sum (x_i-\bar{x}_A)^2 = 1524.90
		\]
		\[
		s_A = \sqrt{\frac{1524.90}{10-1}} = \sqrt{169.43} = 13.02\ \text{USD}
		\]

		\textbf{Sample B}

		Step 1: Compute $\bar{x}_B$.
		\[
		\bar{x}_B = \frac{\sum x_i}{n_B}
		= \frac{54+58+61+63+57+60+56+65+59+62+58}{11}
		= \frac{653}{11}
		= 59.36\ \text{USD}
		\]

		Step 2: Compute $s_B$.
		\[
		s_B = \sqrt{\frac{\sum (x_i-\bar{x}_B)^2}{n_B-1}}
		\]
		Using $\bar{x}_B=59.36$:
		\[
		\sum (x_i-\bar{x}_B)^2 = 104.55
		\]
		\[
		s_B = \sqrt{\frac{104.55}{11-1}} = \sqrt{10.45} = 3.23\ \text{USD}
		\]

		Sample A has much larger spread than Sample B ($13.02$ USD vs $3.23$ USD), so video monetization is more volatile and less reliable for stable monthly cash-flow planning.

		\newpage
		\item
		\subsection*{Problem 3}
		\subsection*{Problem description}
		A teen fintech app tracks referral bonuses paid to student ambassadors in one week. Management wants to know both average bonus size and spread before changing the incentive plan.

		\begin{center}
		\textit{Bonus per student in one week (USD)} \par
		\begin{tabular}{c c}
		\toprule
		Bonus $x$ & Frequency $f$\\
		\midrule
		8  & 3\\
		12 & 5\\
		16 & 7\\
		20 & 4\\
		24 & 2\\
		\bottomrule
		\end{tabular}
		\end{center}

		Treat the full dataset represented by the table as a sample. Compute the sample mean and sample standard deviation. Briefly comment on whether the bonus structure looks tightly clustered or widely dispersed.

		\subsection*{C1}
		We use frequency weighting to compute the sample mean and sample standard deviation.

		Step 1: Compute $\bar{x}$.
		\[
		n = \sum f = 3+5+7+4+2 = 21
		\]
		\[
		\bar{x} = \frac{\sum fx}{n}
		= \frac{8(3)+12(5)+16(7)+20(4)+24(2)}{21}
		= \frac{324}{21}
		= 15.43\ \text{USD}
		\]

		Step 2: Compute $s$.
		\[
		s = \sqrt{\frac{\sum f(x-\bar{x})^2}{n-1}}
		\]
		With $\bar{x}=15.43$:
		\[
		\sum f(x-\bar{x})^2 = 3(8-15.43)^2 + 5(12-15.43)^2 + 7(16-15.43)^2 + 4(20-15.43)^2 + 2(24-15.43)^2 = 457.14
		\]
		\[
		s = \sqrt{\frac{457.14}{21-1}} = \sqrt{22.86} = 4.78\ \text{USD}
		\]

		The average bonus is 15.43 USD, and the spread is about 4.78 USD, so bonuses are moderately clustered around the center rather than extremely wide.

		\newpage
		\item
		\subsection*{Problem 4}
		\subsection*{Problem description}
		A weekend delivery team led by senior students recorded tip income per shift (USD) across a random set of shifts. They need a measure of average tips and variability to forecast whether tips can reliably cover transportation costs.

		\begin{center}
		\textit{Tip income per shift (USD)}\par
		\begin{tabular}{c c}
		\toprule
		Tip amount $x$ & Frequency $f$\\
		\midrule
		6  & 4\\
		9  & 6\\
		12 & 8\\
		15 & 5\\
		18 & 3\\
		\bottomrule
		\end{tabular}
		\end{center}

		Treat this as one sample. Compute the sample mean and sample standard deviation, and state what the spread suggests about uncertainty in shift-level earnings.

		\subsection*{C1}
		We treat all frequency entries as one sample and compute weighted statistics.

		Step 1: Compute $\bar{x}$.
		\[
		n = \sum f = 4+6+8+5+3 = 26
		\]
		\[
		\bar{x} = \frac{\sum fx}{n}
		= \frac{6(4)+9(6)+12(8)+15(5)+18(3)}{26}
		= \frac{303}{26}
		= 11.65\ \text{USD}
		\]

		Step 2: Compute $s$.
		\[
		s = \sqrt{\frac{\sum f(x-\bar{x})^2}{n-1}}
		\]
		Using $\bar{x}=11.65$:
		\[
		\sum f(x-\bar{x})^2 = 347.88
		\]
		\[
		s = \sqrt{\frac{347.88}{26-1}} = \sqrt{13.92} = 3.73\ \text{USD}
		\]

		The mean tip is 11.65 USD per shift, with typical variation of about 3.73 USD, so tip earnings have noticeable but manageable uncertainty.

		\newpage
		\item
		\subsection*{Problem 5}
		\subsection*{Problem description}
		A student-run study-notes subscription channel grouped weekly net earnings (USD) into class intervals to review performance volatility before setting savings targets.

		\begin{center}
		\textit{Weekly net earnings (USD)}\par
		\begin{tabular}{c c}
		\toprule
		Interval & Frequency\\
		\midrule
		40--59   & 3\\
		60--79   & 6\\
		80--99   & 8\\
		100--119 & 5\\
		120--139 & 2\\
		\bottomrule
		\end{tabular}
		\end{center}

		Using class midpoints, approximate the sample mean and sample standard deviation. Explain what your dispersion result implies for setting a realistic weekly savings contribution.

		\subsection*{C1}
		We approximate the grouped sample using class midpoints.

		Step 1: Compute $\bar{x}$ with midpoints.
		\[
		m_i:\ 49.5,\ 69.5,\ 89.5,\ 109.5,\ 129.5
		\qquad
		n = 3+6+8+5+2 = 24
		\]
		\[
		\bar{x} = \frac{\sum f m_i}{n}
		= \frac{49.5(3)+69.5(6)+89.5(8)+109.5(5)+129.5(2)}{24}
		= \frac{2088}{24}
		= 87.00\ \text{USD}
		\]

		Step 2: Compute $s$ with midpoints.
		\[
		s = \sqrt{\frac{\sum f(m_i-\bar{x})^2}{n-1}}
		\]
		\[
		\sum f(m_i-\bar{x})^2 = 12250.00
		\]
		\[
		s = \sqrt{\frac{12250.00}{24-1}} = \sqrt{532.61} = 23.08\ \text{USD}
		\]

		The average weekly earning is about 87.00 USD, while the spread of 23.08 USD shows meaningful volatility, so savings targets should allow for lower-income weeks.

		\newpage
		\item
		\subsection*{Problem 6}
		\subsection*{Problem description}
		A youth investment club recorded one sample of monthly portfolio returns (USD) from low-budget micro-investments. Because returns were stored in grouped form, the club needs midpoint-based estimates to evaluate risk before increasing capital.

		\begin{center}
		\textit{Monthly return intervals (USD)}\par
		\begin{tabular}{c c}
		\toprule
		Interval & Frequency\\
		\midrule
		-20-- -1 & 2\\
		0--19    & 5\\
		20--39   & 7\\
		40--59   & 4\\
		60--79   & 2\\
		\bottomrule
		\end{tabular}
		\end{center}

		Using class midpoints, compute the approximate sample mean and sample standard deviation. Interpret whether the return pattern looks stable enough for conservative student investors.

		\subsection*{C1}
		We estimate center and spread from grouped intervals using class midpoints.

		Step 1: Compute $\bar{x}$.
		\[
		m_i:\ -10.5,\ 9.5,\ 29.5,\ 49.5,\ 69.5
		\qquad
		n = 2+5+7+4+2 = 20
		\]
		\[
		\bar{x} = \frac{\sum f m_i}{n}
		= \frac{(-10.5)(2)+9.5(5)+29.5(7)+49.5(4)+69.5(2)}{20}
		= \frac{570}{20}
		= 28.50\ \text{USD}
		\]

		Step 2: Compute $s$.
		\[
		s = \sqrt{\frac{\sum f(m_i-\bar{x})^2}{n-1}}
		\]
		\[
		\sum f(m_i-\bar{x})^2 = 9980.00
		\]
		\[
		s = \sqrt{\frac{9980.00}{20-1}} = \sqrt{525.26} = 22.92\ \text{USD}
		\]

		The estimated return mean is 28.50 USD, but the spread is almost as large as the mean, so the pattern is not very stable for conservative investors.

		\newpage
		\item
		\subsection*{Problem 7}
		\subsection*{Problem description}
		A student creator sells printable planners online. Most weeks are summarized as discrete grouped sales, but two viral-promotion weeks were outliers and must be included in the same sample for risk analysis.

		\begin{center}
		\textit{Regular-week net sales (USD)}\par
		\begin{tabular}{c c}
		\toprule
		Sales $x$ & Frequency $f$\\
		\midrule
		70  & 4\\
		85  & 6\\
		100 & 5\\
		115 & 3\\
		\bottomrule
		\end{tabular}
		\end{center}

		\textbf{Additional outlier observations (USD):} 185, 210.

		Using all observations together as one sample, compute $\bar{x}$ and $s$. Then explain how the outliers influence measured dispersion and why that matters when budgeting for next month.

		\subsection*{C1}
		We combine the grouped regular weeks with the two outlier observations into one sample.

		Step 1: Compute $\bar{x}$.
		\[
		n = (4+6+5+3)+2 = 20
		\]
		\[
		\bar{x} = \frac{70(4)+85(6)+100(5)+115(3)+185+210}{20}
		= \frac{2030}{20}
		= 101.50\ \text{USD}
		\]

		Step 2: Compute $s$.
		\[
		s = \sqrt{\frac{\sum (x_i-\bar{x})^2}{n-1}}
		= \sqrt{\frac{\sum f(x-\bar{x})^2 + (185-\bar{x})^2 + (210-\bar{x})^2}{19}}
		\]
		\[
		\sum (x_i-\bar{x})^2 = 24905.00
		\]
		\[
		s = \sqrt{\frac{24905.00}{19}} = \sqrt{1310.79} = 36.20\ \text{USD}
		\]

		The outlier weeks pull the mean upward and increase the spread substantially, so budgeting should not rely only on viral-week revenue.

		\newpage
		\item
		\subsection*{Problem 8}
		\subsection*{Problem description}
		A gaming streamer (age 17) grouped normal weekly monetization income (USD) by intervals, but also recorded two sponsorship weeks that are clear outliers. The full set is used to evaluate overall volatility before committing to fixed monthly expenses.

		\begin{center}
		\textit{Regular weekly monetization income (USD)}\par
		\begin{tabular}{c c}
		\toprule
		Interval & Frequency\\
		\midrule
		30--49   & 4\\
		50--69   & 7\\
		70--89   & 6\\
		90--109  & 3\\
		\bottomrule
		\end{tabular}
		\end{center}

		\textbf{Additional outlier observations (USD):} 160, 190.

		Using class midpoints for the grouped part and including the outliers explicitly, compute the sample mean and sample standard deviation. Interpret what the volatility indicates about the risk of depending on streaming as a primary income source.

		\subsection*{C1}
		We use class midpoints for regular weeks, then include both sponsorship outliers in the same sample.

		Step 1: Compute $\bar{x}$.
		\[
		m_i:\ 39.5,\ 59.5,\ 79.5,\ 99.5
		\qquad
		n = (4+7+6+3)+2 = 22
		\]
		\[
		\bar{x} = \frac{39.5(4)+59.5(7)+79.5(6)+99.5(3)+160+190}{22}
		= \frac{1700}{22}
		= 77.27\ \text{USD}
		\]

		Step 2: Compute $s$.
		\[
		s = \sqrt{\frac{\sum (x_i-\bar{x})^2}{n-1}}
		= \sqrt{\frac{\sum f(m_i-\bar{x})^2 + (160-\bar{x})^2 + (190-\bar{x})^2}{21}}
		\]
		\[
		\sum (x_i-\bar{x})^2 = 28981.36
		\]
		\[
		s = \sqrt{\frac{28981.36}{21}} = \sqrt{1380.06} = 37.15\ \text{USD}
		\]

		The spread is large compared with the mean, showing high volatility; this makes streaming income risky as a sole source for fixed monthly commitments.

		\newpage
		\item
		\subsection*{Problem 9}
		\subsection*{Problem description}
		A student-built homework app has three realistic weekend revenue outcomes (USD) based on recent launch analytics. The founder models weekend revenue as a random variable $X$ to plan ad spending.

		\begin{center}
		\textit{Probability distribution of weekend revenue}\par
		\begin{tabular}{c c}
		\toprule
		Outcome $x$ (USD) & Probability $P(X=x)$\\
		\midrule
		40  & 0.25\\
		70  & 0.45\\
		110 & 0.30\\
		\bottomrule
		\end{tabular}
		\end{center}

		Compute $E[X]$, $\mathrm{Var}(X)$, and $\mathrm{SD}(X)$. Interpret the mean as the expected weekend financial outcome and explain how the standard deviation informs revenue risk.

		\subsection*{C1}
		We first compute the weighted mean outcome, then discuss dispersion.

		Step 1: Compute the mean value (expected value).
		\[
		\bar{x}_{\text{weighted}} = \sum x\,P(X=x)
		= 40(0.25)+70(0.45)+110(0.30)
		= 10+31.5+33
		= 74.50\ \text{USD}
		\]
		So $E[X]=74.50$ USD.

		Step 2: Sample standard deviation formula check.
		\[
		s = \sqrt{\frac{\sum (x_i-\bar{x})^2}{n-1}}
		\]
		A sample standard deviation $s$ cannot be determined because no sample size $n$ or sample observations are provided; only a probability model is given.

		For model-based dispersion:
		\[
		\mathrm{Var}(X)=\sum (x-\mu)^2P(X=x)
		= (40-74.5)^2(0.25)+(70-74.5)^2(0.45)+(110-74.5)^2(0.30)
		= 684.75
		\]
		\[
		\mathrm{SD}(X)=\sqrt{684.75}=26.17\ \text{USD}
		\]

		The model predicts an average weekend revenue of 74.50 USD, with a typical uncertainty of about 26.17 USD around that expected value.

		\newpage
		\item
		\subsection*{Problem 10}
		\subsection*{Problem description}
		A student startup sells a mini online course and records net revenue (USD) from 40 independent ad-campaign launches. The team wants to model campaign outcome as a random variable by converting observed frequencies into probabilities.

		\begin{center}
		\textit{Net revenue by campaign outcome class}\par
		\begin{tabular}{c c}
		\toprule
		Revenue outcome $x$ (USD) & Frequency\\
		\midrule
		20  & 8\\
		50  & 12\\
		80  & 10\\
		120 & 6\\
		150 & 4\\
		\bottomrule
		\end{tabular}
		\end{center}

		First compute the probability distribution from the grouped counts. Then compute $\mu=E[X]$, $\mathrm{Var}(X)$, and $\mathrm{SD}(X)$. In context, explain why these statistics are useful before scaling campaign spending.

		\subsection*{C1}
		We convert the frequencies to probabilities, then compute the mean and spread.

		Step 1: Compute $\bar{x}$ (and $E[X]$ here).
		\[
		n=8+12+10+6+4=40
		\]
		\[
		P(X=20)=\frac{8}{40}=0.20,\quad
		P(X=50)=\frac{12}{40}=0.30,\quad
		P(X=80)=\frac{10}{40}=0.25,
		\]
		\[
		P(X=120)=\frac{6}{40}=0.15,\quad
		P(X=150)=\frac{4}{40}=0.10
		\]
		\[
		\bar{x}=\frac{\sum fx}{n}
		=\frac{20(8)+50(12)+80(10)+120(6)+150(4)}{40}
		=\frac{2880}{40}
		=72.00\ \text{USD}
		\]
		Thus $\mu=E[X]=72.00$ USD.

		Step 2: Compute $s$ from the sample and also report model spread.
		\[
		s=\sqrt{\frac{\sum f(x-\bar{x})^2}{n-1}}
		\]
		\[
		\sum f(x-\bar{x})^2 = 66240.00
		\]
		\[
		s=\sqrt{\frac{66240.00}{39}}=\sqrt{1698.46}=41.21\ \text{USD}
		\]

		For the random-variable model,
		\[
		\mathrm{Var}(X)=\sum (x-\mu)^2P(X=x)=1656.00
		\]
		\[
		\mathrm{SD}(X)=\sqrt{1656.00}=40.69\ \text{USD}
		\]

		The expected campaign revenue is 72.00 USD, and the sizeable spread (about 41 USD) shows that outcomes vary substantially, so scaling ad spending requires careful risk control.

	\end{ExamProblems}
\end{document}
